
\begin{abstract}
% Large language models (LLMs) are rapidly evolving from conversational systems into sophisticated reasoners capable of tackling complex tasks such as Olympiad-level mathematics and competitive programming. While scaling model size and test-time computation has driven much of this progress, a persistent bottleneck remains: the scarcity of high-quality training problems. Human-curated datasets are costly and limited in scope, while existing synthetic corpora often lack sufficient difficulty and diversity, hindering their effectiveness for post-training strong models. PromptCoT~1.0 recently introduced the rationale-based paradigm for prompt synthesis, showing that incorporating intermediate reasoning steps can significantly increase the difficulty of generated problems. Building on this insight, we present PromptCoT~2.0, a principled and scalable framework that moves beyond hand-crafted instructions and domain-specific heuristics. Our method introduces an expectation–maximization (EM) optimization loop, where rationales are iteratively refined to guide prompt construction, yielding problems that are substantially harder and more varied than those in existing open-source corpora. The resulting synthetic prompts serve as a foundation for two complementary post-training regimes: (1) \emph{self-play}, where strong models improve autonomously from verifiable feedback without requiring stronger external teachers; and (2) \emph{supervised fine-tuning (SFT)}, where weaker models are trained on reasoning traces distilled from a teacher. Extensive experiments demonstrate the effectiveness of this approach. In the self-play setting, applying PromptCoT~2.0 to \texttt{Qwen3-30B-A3B-Thinking-2507} establishes new state-of-the-art results \emph{at the 30B scale}, with accuracy gains of +4.4, +4.8, and +5.3 points on AIME~24/25 and HMMT~Feb~25, +6.1 and +5.0 on LiveCodeBench v5/v6, and +35 Elo on Codeforces.  In the SFT setting, training \texttt{Qwen2.5-7B-Instruct} solely on our synthetic prompts boosts performance from weak baselines to strong, competitive results across all six benchmarks (e.g., 73.1 on AIME~24, 65.6 on AIME~25, and 53.4 on LiveCodeBench v5), demonstrating that purely synthesized problems can achieve performance on par with human-curated or hybrid datasets. Beyond accuracy, distributional and difficulty analyses confirm that PromptCoT~2.0 produces problems that differ fundamentally from prior datasets and demand substantially deeper reasoning. These results establish prompt synthesis as a powerful new axis for scaling LLM reasoning, and position PromptCoT~2.0 as a scalable foundation for the next generation of open-source reasoning models. The implementation is available at \url{https://github.com/inclusionAI/PromptCoT}.

Large language models (LLMs) are evolving from conversational systems into strong reasoners for tasks such as Olympiad mathematics and competitive programming. While scaling parameters and test-time computation has driven progress, a key bottleneck is the lack of high-quality training problems: human-curated datasets are costly and limited, while existing synthetic corpora are often too easy or narrow. PromptCoT 1.0 showed that injecting rationales into prompt synthesis increases problem difficulty. Building on this, we present PromptCoT 2.0, a scalable framework that replaces hand-crafted heuristics with an expectation–maximization (EM) loop, where rationales are iteratively refined to guide prompt construction. This produces problems that are both harder and more diverse than prior corpora. The synthetic prompts support two post-training regimes: (1) \emph{Self-Play}, where strong models improve autonomously via verifiable feedback without stronger teachers; and (2) \emph{Supervised Fine-Tuning (SFT)}, where weaker models learn from teacher-distilled traces. Extensive experiments demonstrate the effectiveness of this approach. In self-play, applying PromptCoT 2.0 to \texttt{Qwen3-30B-A3B-Thinking-2507} sets new state-of-the-art results \emph{at the 30B scale}, with +4.4, +4.8, and +5.3 on AIME 24/25 and HMMT 25, +6.1 and +5.0 on LiveCodeBench v5/v6, and +35 Elo on Codeforces. In SFT, training \texttt{Qwen2.5-7B-Instruct} solely on synthetic prompts boosts accuracy to 73.1 (AIME 24), 65.6 (AIME 25), and 53.4 (LiveCodeBench v5), surpassing models trained on human or hybrid data. Analyses further confirm that PromptCoT 2.0 yields fundamentally harder and distributionally distinct problems. These results establish prompt synthesis as a new axis for scaling reasoning and position PromptCoT 2.0 as a scalable foundation for future open-source models. The implementation is available at \url{https://github.com/inclusionAI/PromptCoT}.
\end{abstract}

